% ---
% title: The recovery rate
% tags: finance, density, distribution-function, expectation, variance
% ---

\begin{problem}
	\begin{context}
		When a borrower defaults, the bank rarely loses the whole loan: it sells
		the collateral, or it gets part of the money back in a court procedure.
		The \emph{recovery rate} is the fraction of the defaulted loan that the
		bank gets back. It is a number between 0 (nothing is recovered) and 1
		(everything is recovered), and before the default it is not known.
		Banks need its expected value and its spread to set aside enough
		capital for the losses on their loans.
	\end{context}
	The recovery rate $R$ of a defaulted loan is a continuous random variable
	with density
	\[
		f(x)=\begin{cases}
			6x(1-x) & 0\le x\le 1, \\
			0       & \text{otherwise}.
		\end{cases}
	\]
	(a) Find the distribution function $F$ of $R$.
	(b) Find $P(R>0.8)$.
	(c) Find $P(0.2<R<0.4)$.
	(d) Find $\E[R]$, $\Var(R)$ and $\SD(R)$.
	(e) A loan of 1 million euros has defaulted, and the bank loses
	$1\,000\,000\,(1-R)$ euros on it. Find the expected loss and the standard
	deviation of the loss.
\end{problem}

\begin{slideproblem}
	The recovery rate $R$ (the fraction of a defaulted loan that the bank
	gets back) has density $f(x)=6x(1-x)$ for $0\le x\le1$, and $0$ otherwise.
	(a) Find the distribution function $F$.
	(b) $P(R>0.8)$. \ (c) $P(0.2<R<0.4)$.
	(d) $\E[R]$, $\Var(R)$, $\SD(R)$.
	(e) On a defaulted loan of 1 million euros the bank loses
	$1\,000\,000\,(1-R)$ euros: expected loss and its standard deviation?
\end{slideproblem}

\begin{slides}
	\small
	\exstatement
	\pause
	\begin{columns}[T]
		\begin{column}{.47\textwidth}
			\sol (a) For $0\le x\le1$:\\
			$F(x)=\int_0^x(6t-6t^2)\,dt=\alert{3x^2-2x^3}$;\\
			$F(x)=0$ for $x<0$, \ $F(x)=1$ for $x>1$.\\[1mm]
			\pause
			(b) $P(R>0.8)=1-F(0.8)=\alert{0.104}$\\[1mm]
			(c) $F(0.4)-F(0.2)=0.352-0.104=\alert{0.248}$
		\end{column}
		\begin{column}{.51\textwidth}
			\pause
			(d) $\E[R]=\int_0^1x\cdot6x(1-x)\,dx=2-\frac32=\alert{\frac12}$\\[1mm]
			$\E[R^2]=\int_0^1x^2\cdot6x(1-x)\,dx=\frac32-\frac65=\frac3{10}$\\[1mm]
			$\Var(R)=\frac3{10}-\frac14=\alert{\frac1{20}}$, \
			$\SD(R)=\sqrt{0.05}\approx\alert{0.2236}$\\[1mm]
			\pause
			(e) $\E[\text{loss}]=1\,000\,000\,(1-\E[R])=\alert{500\,000}$ euros\\[1mm]
			$\SD(\text{loss})=1\,000\,000\cdot\SD(R)\approx\alert{223\,607}$ euros
			\ (as $\Var(aX+b)=a^2\Var(X)$)
		\end{column}
	\end{columns}
\end{slides}

\begin{solution}
	\idea The distribution function is the area under the density to the left
	of $x$: $F(x)=P(R\le x)=\int_{-\infty}^{x}f(t)\,dt$. Once $F$ is known,
	every interval probability is a difference of two values of $F$. The
	expected value and the variance are integrals of $x\,f(x)$ and
	$x^2f(x)$; the loss is a linear function of $R$, so its expected value
	and standard deviation follow from those of $R$.

	\step{The picture.} The density is a hill on $[0,1]$: it is $0$ at both
	ends and largest in the middle, $f(0.5)=6\cdot0.5\cdot0.5=1.5$, and it is
	symmetric about $0.5$. Each
	probability below is an area under this curve. The shaded parts show
	$P(0.2<R<0.4)$ (teal) and $P(R>0.8)$ (gold).
	\begin{center}
		\begin{tikzpicture}[x=6cm,y=2.4cm,>=stealth]
			\fill[teal!35] (0.2,0)
			-- plot[domain=0.2:0.4,samples=30] (\x,{6*\x*(1-\x)}) -- (0.4,0) -- cycle;
			\fill[gold!60] (0.8,0)
			-- plot[domain=0.8:1,samples=30] (\x,{6*\x*(1-\x)}) -- (1,0) -- cycle;
			\draw[->] (-0.1,0)--(1.15,0) node[right] {$x$};
			\draw[->] (0,-0.05)--(0,1.75) node[above] {$f(x)$};
			\foreach \x in {0.2,0.4,0.5,0.8,1}
			\draw (\x,0.03)--(\x,-0.03) node[below] {\footnotesize$\x$};
			\node[below left] at (0,0) {\footnotesize$0$};
			\draw[gray!60,dashed] (0,1.5)--(0.5,1.5);
			\node[left] at (0,1.5) {\footnotesize$1.5$};
			\draw[navy,very thick] (-0.1,0)--(0,0);
			\draw[navy,very thick] plot[domain=0:1,samples=60] (\x,{6*\x*(1-\x)});
			\draw[navy,very thick] (1,0)--(1.12,0);
		\end{tikzpicture}
	\end{center}
	Before we start, check that $f$ is a density: it is never negative, and
	the total area is
	\[
		\int_0^1 (6x-6x^2)\,dx=\Bigl[3x^2-2x^3\Bigr]_0^1=3-2=1.
	\]

	\step{(a)} We compute $F(x)=\int_{-\infty}^{x}f(t)\,dt$ separately on the
	three pieces where $f$ has a different formula.
	\begin{itemize}
		\item For $x<0$: the density is $0$ to the left of $x$, so there is no
		      area, and $F(x)=0$. (The recovery rate cannot be negative.)
		\item For $0\le x\le1$: the density is $0$ left of $0$, so only the
		      part from $0$ to $x$ counts. Write the density as $6t-6t^2$. An
		      antiderivative is $3t^2-2t^3$ (check: its derivative is
		      $6t-6t^2$). So
		      \[
			      F(x)=\int_0^x (6t-6t^2)\,dt=\Bigl[3t^2-2t^3\Bigr]_0^x
			      =(3x^2-2x^3)-(0-0)=3x^2-2x^3.
		      \]
		\item For $x>1$: all the area (which is $1$) lies to the left of $x$,
		      so $F(x)=1$. Indeed, the formula above gives $F(1)=3-2=1$, and
		      the density adds nothing after $1$.
	\end{itemize}
	So
	\[
		F(x)=\begin{cases}
			0          & x<0,          \\
			3x^2-2x^3  & 0\le x\le 1,  \\
			1          & x>1.
		\end{cases}
	\]
	Check: $F'(x)=6x-6x^2=f(x)$ on $(0,1)$, as it must be.

	\step{(b)} ``$R>0.8$'' is the complement of ``$R\le0.8$'', so
	$P(R>0.8)=1-F(0.8)$. Here
	\[
		F(0.8)=3\cdot0.64-2\cdot0.512=1.92-1.024=0.896,
	\]
	so
	\[
		P(R>0.8)=1-0.896=0.104=\frac{13}{125}.
	\]
	The gold area in the picture is small: near $1$ the density is low, so a
	near-complete recovery is rare.

	\step{(c)} Use $P(a<R\le b)=F(b)-F(a)$. $R$ is continuous, so a single
	point has probability $0$: $P(R=0.4)=0$, and the strict inequality $R<0.4$
	changes nothing. First the two values of $F$:
	\[
		F(0.4)=3\cdot0.16-2\cdot0.064=0.48-0.128=0.352,
	\]
	\[
		F(0.2)=3\cdot0.04-2\cdot0.008=0.12-0.016=0.104.
	\]
	Then
	\[
		P(0.2<R<0.4)=F(0.4)-F(0.2)=0.352-0.104=0.248=\frac{31}{125}.
	\]
	(Notice $F(0.2)=0.104=1-F(0.8)$: the density is symmetric about $0.5$,
	$f(1-x)=6(1-x)x=f(x)$, so the area left of $0.2$ equals the area right of
	$0.8$.)

	\step{(d)} The expected value of a continuous random variable is
	$\E[R]=\int_{-\infty}^{\infty}x\,f(x)\,dx$, the analogue of
	$\sum x\,p(x)$ in the discrete case. Outside $[0,1]$ the density is $0$:
	\[
		\E[R]=\int_0^1 x\cdot6x(1-x)\,dx=\int_0^1(6x^2-6x^3)\,dx
		=\Bigl[2x^3-\frac{3x^4}{2}\Bigr]_0^1=2-\frac32=\frac12 .
	\]
	This was to be expected: the density is symmetric about $0.5$, so its
	balance point is $0.5$. For the variance we use
	$\Var(R)=\E[R^2]-(\E[R])^2$, and $\E[R^2]=\int x^2f(x)\,dx$:
	\[
		\E[R^2]=\int_0^1 x^2\cdot6x(1-x)\,dx=\int_0^1(6x^3-6x^4)\,dx
		=\Bigl[\frac{3x^4}{2}-\frac{6x^5}{5}\Bigr]_0^1=\frac32-\frac65=\frac3{10}.
	\]
	So
	\[
		\Var(R)=\frac3{10}-\Bigl(\frac12\Bigr)^2=\frac3{10}-\frac14=\frac1{20}=0.05,
		\qquad
		\SD(R)=\sqrt{0.05}\approx0.2236 .
	\]
	On average the bank recovers half of a defaulted loan, and the recovery
	rate typically differs from $0.5$ by about $0.22$.

	\step{(e)} The loss is $1\,000\,000\,(1-R)=-1\,000\,000\,R+1\,000\,000$
	euros, a linear function $aR+b$ of $R$ with $a=-1\,000\,000$ and
	$b=1\,000\,000$. For a linear function, $\E[aR+b]=a\,\E[R]+b$ and
	$\Var(aR+b)=a^2\Var(R)$, so $\SD(aR+b)=|a|\,\SD(R)$ (adding a constant does
	not change the spread, and the sign of $a$ does not matter). Hence
	\[
		\E[\text{loss}]=1\,000\,000\,\bigl(1-\E[R]\bigr)=1\,000\,000\cdot\frac12
		=500\,000 \text{ euros},
	\]
	\[
		\SD(\text{loss})=1\,000\,000\cdot\SD(R)=1\,000\,000\cdot\sqrt{0.05}
		\approx223\,607 \text{ euros}.
	\]
	The bank expects to lose half of the loan, but the actual loss is quite
	uncertain: a typical deviation from $500\,000$ euros is more than
	$200\,000$ euros.

	\begin{answer}
		(a) $F(x)=0$ for $x<0$, $F(x)=3x^2-2x^3$ for $0\le x\le1$, $F(x)=1$
		for $x>1$ \quad
		(b) $P(R>0.8)=0.104$ \quad
		(c) $P(0.2<R<0.4)=0.248$ \quad
		(d) $\E[R]=\frac12$, $\Var(R)=\frac1{20}$, $\SD(R)\approx0.2236$ \quad
		(e) expected loss $500\,000$ euros, standard deviation
		$\approx223\,607$ euros
	\end{answer}
	\begin{remark}
		\textbf{Common mistake:} reading a probability off the density,
		e.g.\ saying $P(R=0.5)=f(0.5)=1.5$. A density value is not a
		probability (here it is even larger than $1$). Probabilities are
		\emph{areas} under $f$, and the area above a single point is $0$.
		Another one: computing the variance in (e) as $1\,000\,000\cdot\Var(R)$.
		The factor comes in squared, $\Var(aR+b)=a^2\Var(R)$; only the standard
		deviation scales with $|a|$.

		The loss given default is $L=1-R$, the fraction of the loan that is
		lost (compare Exercise~\exref{loss-given-default}). Since the density
		of $R$ is symmetric about $0.5$, here $L$ has the same distribution
		as $R$.
	\end{remark}
\end{solution}
